\section{Related Work}

\noindent\textbf{Segment Anything Model.}
SAM~\cite{kirillov2023sam} introduced promptable segmentation trained on the SA-1B dataset. SAM2~\cite{ravi2024sam2} extended this to video. Several works adapt SAM to specialized domains: SAM-Adapter~\cite{chen2023samadapter} adds learnable modules \emph{inside} the image encoder, while COMPrompter~\cite{hao2024comprompter} reconceptualizes SAM with multi-prompt networks. Liu~\etal~\cite{liu2025dualstream} proposed dual-stream adapters. Our approach differs in two ways: we fine-tune the \emph{existing} decoder without adding parameters, and our LLPM operates \emph{before} the encoder in input space---a placement unexplored by prior work.

\noindent\textbf{Camouflaged Object Detection.}
Dedicated COD architectures include SINet~\cite{fan2020cod10k} (search-and-identify), PFNet~\cite{mei2021pfnet} (positioning and focus), ZoomNet~\cite{pang2022zoomnet} (mixed-scale learning), and SegMaR~\cite{jia2022segmar}. These design specialized networks from scratch. We instead leverage SAM's pre-trained representations and show that decoder-only adaptation achieves competitive improvements. Das~\etal~\cite{das2025camouflage} recently introduced synthetic camouflage generation, which is complementary to our detection-side approach.

\noindent\textbf{Parameter-Efficient Fine-Tuning.}
LoRA~\cite{hu2022lora} and prompt tuning~\cite{lester2021prompt} are standard for adapting large models. For SAM, these techniques add parameters to the encoder or prompt pathway. Our decoder-only approach is simpler: we freeze everything except the existing mask decoder, adding zero new parameters. The LLPM extends this with a minimal 31K-parameter pre-encoder module.
